\begin{abstract}	
\noindent 
This paper investigates the persistent debate in the literature on whether informal employment is a choice driven by comparative advantages and off-the-job benefits, or a necessity due to barriers to formal employment. To address this, we consider the unobservable characteristics influencing workers' sector choices and earnings, aiming to distinguish between labor market segmentation and comparative advantage. The study's main objective is to examine whether an occupation-based wage premium exists for formal employment compared to informal employment across the wage distribution in a developing economy. Using quantile regression, we capture the heterogeneity in workers' responses to observed attributes and measure the wage gap at different quantiles of the wage distribution. We apply the unconditional quantile regression method and RIF-OLS decomposition techniques developed by Firpo and Nopo. Our results reveal a significant wage gap between formal and informal employment, with the smallest gap among elementary workers, increasing disparity for clerical workers, and the largest gap at the median for managers across both rounds.\par
\vspace{.1in}
\noindent\textbf{Keywords:} Labor market wage discrimination, Heterogeneity, Job occupation informality, RIF decomposition, Unconditional quantile regression \\
\noindent\textbf{JEL Codes:} J31, J24, J46, C21, C14
\bigskip
\end{abstract}